\scanpage{83}
\diagram{\begin{tikzcd}[ampersand replacement=\&,column sep=small]
\&\&\&\underline P_{X/S}\arrow[dl,"{\unsure{\hat d|\dot d|j}}",swap]\arrow[d]\\
0\arrow[r,dash]\&\underline{\omega}_{X/S}\arrow[r]\&f_*(\underline E_{X/S})\arrow[r]\&f_*(\underline{\mathrm{Div}}_{X/S})^{w}\arrow[r,dash]\&0
\end{tikzcd}}
définie par
\[
f_*(\underline R^{\unsure{\times|*}}_{X/S})\to f_*(E_{\deleted{\unsure{Y}}X/S})
\]
donnée par
\[
f\mapsto\deleted{\unsure{f}}\left(\mathrm{div}\,f,\frac{df}{f}\right)
\]

Il faut vérifier que c'est nul sur $f_*(\underline O^*_X)$, cela marche si \inserted{$f$ \uuline{propre et plat}, \deleted{$\underline O_X$} $\underline O_X$ \uuline{est \unsure{coh.} plat en dim $0$ sur $S$ et}} $f_*(\underline O_X)$ \uuline{est étale sur $S$} (p.ex. si $\mathcal O_S\simeq f_*(\underline O_X)$ universellement), ce que nous allons supposer.
\margin{En l'absence de cette hypothèse, on obtient seulement une ext. de $\underline{\mathrm{Pic}}^w_{X/S}$ par $\underline{\omega}_{X/S}/\operatorname{Im} f_*(\underline O_X^*)$}
\diagram{\begin{tikzcd}[ampersand replacement=\&,row sep=tiny,column sep=tiny]
\&X\arrow[dl]\\ S'\arrow[dr]\\ \&S
\end{tikzcd}}
Nous obtenons alors une extension
\[
\mathcal M_{X/S}=f_*(\underline E_{\deleted{\unsure{Y}}X/S})/\underline P_{X/S}
\]
\diagram{\begin{tikzcd}[ampersand replacement=\&,column sep=small]
0\arrow[r]\&\underline{\omega}_{X/S}\arrow[r]\arrow[d,dash]\&\mathcal M_{X/S}\arrow[r]\&\underline{\mathrm{Pic}}^w_{X/S}\arrow[r]\&0\\
\&\substack{\text{fibré vectoriel}\\ \text{sur }S}
\end{tikzcd}}

\uline{Cas favorable} $\underline{\Omega}^1_{X/S}$ \uuline{est \unsure{ch.|coh.} plat sur} $S$ \uuline{en} dim $0$, i.e. $\underline{\omega}_{X/S}$ est un fibré vectoriel associé à un module \uline{loc. libre} sur $S$. [C'est le cas si $f$ est lisse et projectif et $S$ de car nulle (?), ou si $X/S$ est un schéma abélien]. Alors $R^1 f_*(\underline{\Omega}^1_{X/S})$ est aussi représentable par un fibré vectoriel, donc si $\underline{\mathrm{Pic}}_{X/S}$ est représentable (p.ex. $f$ à fibres \unsure{géom.|géom} intègres) $\underline{\mathrm{Pic}}^w_{X/S}$ l'est aussi. Donc par descente $\mathcal M_{X/S}$ l'est également.
% archive sheet number: 82

